\documentclass[11pt,a4paper]{article}

\usepackage[utf8]{inputenc}
\usepackage[T1]{fontenc}
\usepackage{lmodern}
\usepackage[margin=2.5cm]{geometry}
\usepackage{amsmath,amssymb,amsthm}
\usepackage{booktabs}
\usepackage{graphicx}
\usepackage{hyperref}
\usepackage{algorithm}
\usepackage{algorithmic}
\usepackage{natbib}
\usepackage{xcolor}
\usepackage{tikz}
\usetikzlibrary{arrows.meta,positioning,shapes.geometric,fit}

\hypersetup{
    colorlinks=true,
    linkcolor=blue!60!black,
    citecolor=green!50!black,
    urlcolor=blue!70!black,
}

\newcommand{\RR}{\mathbb{R}}
\newcommand{\bx}{\mathbf{x}}
\newcommand{\bW}{\mathbf{W}}
\newcommand{\bP}{\mathbf{P}}
\newcommand{\bE}{\mathbf{E}}
\newcommand{\bI}{\mathbf{I}}

\title{EigenDialectos: Spectral Dialect Classification\\for Spanish Varieties via Contrastive Learning\\and Eigenmode Fingerprinting}

\author{Jose Luis Saorin}
\date{April 2026}

\begin{document}
\maketitle

\begin{abstract}
We present EigenDialectos, a system for classifying Spanish text into eight
regional dialect varieties using two complementary approaches built on a
shared transformer backbone. The first method fine-tunes BETO (a Spanish BERT
model) with LoRA adapters using a multi-task objective combining masked language
modeling, ArcFace angular-margin classification, and supervised contrastive
learning with Momentum Contrast (MoCo) queues. The second method extracts
contextual word embeddings per variety, aligns them via Procrustes, computes
inter-dialect transformation matrices, and classifies text via eigenmode
spectral fingerprints. On a test set of dialectally-marked sentences, the
projection-space centroid classifier achieves 58\% accuracy across 8 varieties
(chance = 12.5\%), with strong performance on lexically distinctive dialects
(Rioplatense, Chilean, Mexican) and room for improvement on varieties with
shared Iberian features (Andalusian, Canarian). We describe the architecture,
training dynamics, the ArcFace collapse phenomenon and its resolution, and the
mathematical framework underlying the spectral approach.
\end{abstract}

\section{Introduction}

Spanish is spoken by over 500 million people across more than 20 countries,
giving rise to significant regional variation at lexical, phonological,
morphosyntactic, and pragmatic levels. While dialect identification has been
studied for major world languages \citep{zampieri2017findings}, fine-grained
classification of Spanish varieties---particularly distinguishing between
closely related dialects like Andalusian, Canarian, and Peninsular
standard---remains challenging.

We address this with EigenDialectos, which combines modern contrastive
representation learning with classical spectral analysis of linear
transformations. Our contributions are:

\begin{enumerate}
    \item A multi-task training framework combining MLM, ArcFace, and
    SupCon+MoCo+DCL contrastive objectives for dialect-discriminative
    representations.
    \item A spectral classification method based on eigendecomposition of
    inter-dialect transformation matrices.
    \item Analysis of the ArcFace angular collapse phenomenon when
    classification and contrastive objectives operate in misaligned feature
    spaces.
    \item A projection-space centroid classifier that exploits the learned
    contrastive geometry for zero-shot dialect identification.
\end{enumerate}

\section{Related Work}

\paragraph{Dialect identification.}
The VarDial shared tasks \citep{zampieri2017findings} established benchmarks
for discriminating between language varieties. For Spanish specifically,
\citet{rangel2016overview} addressed author profiling across varieties.
Most approaches use n-gram features or fine-tuned transformers with
cross-entropy classification.

\paragraph{Contrastive learning for NLP.}
Supervised contrastive learning \citep{khosla2020supervised} extends
self-supervised contrastive methods to leverage label information. The
Decoupled Contrastive Learning (DCL) objective \citep{yeh2022decoupled}
improves gradient quality by removing positives from the denominator.
MoCo \citep{he2020momentum} maintains a momentum-updated encoder to
produce consistent negative examples via a queue, avoiding the staleness
problem of cross-batch memory.

\paragraph{Angular margin losses.}
ArcFace \citep{deng2019arcface} adds an additive angular margin to the target
class logit in the cosine classification space, encouraging large angular
separation between classes. It has been widely used in face recognition and
speaker verification.

\section{Model Architecture}

\subsection{Base Encoder}

We use BETO \citep{canete2020spanish}, a Spanish BERT model pre-trained on a
large Spanish corpus with 110M parameters. All base parameters are frozen. We
apply Low-Rank Adaptation (LoRA) \citep{hu2022lora} with rank $r=16$ and
$\alpha=32$ to the attention query, key, value, and output projections, as well
as the feed-forward intermediate and output dense layers across all 12
transformer layers. This yields approximately 2.6M trainable parameters in the
backbone.

\subsection{Variety Tokens}

We add 8 special tokens \texttt{[VAR\_ES\_PEN]}, \ldots,
\texttt{[VAR\_ES\_AND\_BO]} to the vocabulary. During training, the
corresponding variety token is prepended to each input text, providing an
explicit dialect signal to the model. The token embeddings are initialized from
the existing embedding distribution and learned end-to-end.

\subsection{Dialect-Aware Attention Pooling}

Instead of using only the \texttt{[CLS]} token representation, we learn a
weighted average over all token positions:

\begin{equation}
    \mathbf{h}_{\text{pool}} = \sum_{i=1}^{L} \alpha_i \mathbf{h}_i, \quad
    \alpha_i = \frac{\exp(\mathbf{w}^\top \mathbf{h}_i)}
    {\sum_{j=1}^{L} \exp(\mathbf{w}^\top \mathbf{h}_j)}
\end{equation}

where $\mathbf{w} \in \RR^{768}$ is a learned query vector and the sum runs
over non-padding positions. This allows the model to attend to
dialect-significant tokens (regionalisms, morphological markers) rather than
relying solely on the \texttt{[CLS]} aggregation.

\subsection{Projection Head}

The pooled representation $\mathbf{h}_{\text{pool}} \in \RR^{768}$ is passed
through a two-layer MLP with BatchNorm:

\begin{equation}
    \mathbf{z} = \text{Linear}_{384 \to 384}\Big(
    \text{ReLU}\big(\text{BN}\big(\text{Linear}_{768 \to 384}
    (\mathbf{h}_{\text{pool}})\big)\big)\Big)
\end{equation}

The contrastive loss operates on the $\ell_2$-normalized version
$\hat{\mathbf{z}} = \mathbf{z} / \|\mathbf{z}\|$.

\subsection{ArcFace Classifier}

We apply ArcFace angular-margin classification on the \emph{projected}
features $\mathbf{z} \in \RR^{384}$ (before $\ell_2$-normalization):

\begin{equation}
    \text{logit}_j = s \cdot \cos(\theta_j + m \cdot \mathbf{1}[j = y])
    \label{eq:arcface}
\end{equation}

where $\theta_j = \arccos(\hat{\mathbf{z}}^\top \hat{\mathbf{w}}_j)$ is the
angle between the normalized input and the $j$-th class weight vector,
$s = 30$ is the scale factor, $m = 0.3$ is the angular margin, and $y$ is
the true variety label.

\paragraph{ArcFace collapse and resolution.}
In our initial architecture, ArcFace operated on the raw pooled representation
$\mathbf{h}_{\text{pool}} \in \RR^{768}$, while contrastive learning operated
on the projected $\hat{\mathbf{z}} \in \RR^{384}$. This created a
\emph{feature space misalignment}: the projection head learned to extract
dialect-discriminative features from $\mathbf{h}_{\text{pool}}$, but the raw
pooled vector itself did not develop good angular structure for direct cosine
classification. At inference, all ArcFace logits collapsed to $-30.0$
(cosine $\approx -1.0$ with all class weight vectors), indicating that
$\mathbf{h}_{\text{pool}}$ had drifted to a region antipodal to all class
weights.

The fix is straightforward: route ArcFace through the projection space, so
both objectives share the same feature geometry. This ensures that the angular
separation enforced by ArcFace directly benefits the contrastive space, and
vice versa.

\section{Training}

\subsection{Multi-Task Objective}

The total loss combines four terms:

\begin{equation}
    \mathcal{L} = w_{\text{mlm}} \mathcal{L}_{\text{mlm}}
    + w_{\text{cls}} \mathcal{L}_{\text{cls}}
    + w_{\text{con}} \mathcal{L}_{\text{con}}
    + w_{\text{ctr}} \mathcal{L}_{\text{ctr}}
\end{equation}

\paragraph{MLM loss.}
Standard masked language modeling with 15\% masking probability, using BETO's
word embeddings as the prediction head (weight-tied).

\paragraph{CLS loss.}
Cross-entropy on ArcFace logits (Eq.~\ref{eq:arcface}).

\paragraph{Contrastive loss.}
Supervised contrastive with the DCL variant:

\begin{equation}
    \mathcal{L}_{\text{con}} = -\frac{1}{|P_i|} \sum_{j \in P_i}
    \log \frac{\exp(\hat{\mathbf{z}}_i^\top \hat{\mathbf{z}}_j / \tau)}
    {\sum_{k \in N_i} \exp(\hat{\mathbf{z}}_i^\top \hat{\mathbf{z}}_k / \tau)}
    \label{eq:supcon}
\end{equation}

where $P_i$ is the set of positives (same variety, excluding self), $N_i$ is
the set of negatives (different variety), and $\tau = 0.07$ is the temperature.
The DCL formulation \citep{yeh2022decoupled} removes positives from the
denominator, improving gradient quality.

\paragraph{Center loss.}
During the first 2000 contrastive optimization steps, we add center loss
\citep{wen2016discriminative} to pull embeddings toward learned class centroids,
breaking initial symmetry.

\subsection{Curriculum}

Training proceeds in two phases:

\begin{enumerate}
    \item \textbf{Pretrain (epochs 1--2):} MLM + CLS only
    ($w_{\text{mlm}} = 0.5$, $w_{\text{cls}} = 0.5$, $w_{\text{con}} = 0$).
    \item \textbf{Full (epochs 3--10):} All losses active, with curriculum
    shift from $w_{\text{mlm}} = 0.3 \to 0.15$ and
    $w_{\text{con}} = 0.4 \to 0.55$.
\end{enumerate}

\subsection{Momentum Contrast (MoCo)}

With balanced batches of $k = 4$ samples per variety ($B = 32$), in-batch
contrastive has only 24 negatives per query---insufficient for the low
temperature $\tau = 0.07$. We augment with a MoCo momentum queue
\citep{he2020momentum}.

\paragraph{Momentum encoder.}
A deep copy of the full model (including LoRA adapters) is maintained as an
exponential moving average (EMA):

\begin{equation}
    \theta_{\text{mom}} \leftarrow m \cdot \theta_{\text{mom}}
    + (1 - m) \cdot \theta_{\text{train}}
\end{equation}

with momentum $m = 0.999$. Only trainable parameters (LoRA, attention pool,
projection head) are EMA-updated; frozen BETO weights are identical in both
encoders. BatchNorm running statistics are also EMA-updated.

\paragraph{Queue.}
A FIFO circular buffer of size $Q = 4096$ stores $(\hat{\mathbf{z}}_{\text{mom}}, y)$
pairs from the momentum encoder. The queue activates at epoch 4 and ramps
linearly from 128 to 4096 entries over 5000 optimizer steps.

\paragraph{Why not cross-batch memory (XBM)?}
Our initial implementation used XBM \citep{wang2020cross}, which stores
embeddings from the training encoder directly. These go stale as the model
updates. With $\tau = 0.07$, a cosine staleness of $\delta = 0.1$ creates
$\exp(\delta / \tau) \approx 4.2\times$ error in softmax weights. The
contrastive loss plateaued at $\ln(4096) \approx 8.3$ (random baseline).
MoCo solves this because the momentum encoder changes slowly, making all
queue entries mutually consistent.

\subsection{Optimization}

We use AdamW with $\text{lr} = 2 \times 10^{-4}$, linear warmup over 5000
steps, and cosine decay. The projection head receives $10\times$ the base
learning rate. Gradient clipping at norm 1.0.

\section{Spectral Eigenmode Classification}

Beyond the direct neural classifier, we develop a spectral approach based on
eigendecomposition of inter-dialect transformation matrices.

\subsection{Contextual Word Embedding Extraction}

For a vocabulary $V$ of $|V|$ words, we extract contextual embeddings per
variety. For word $w$ in variety $v$, we collect all occurrences of $w$ in
the variety-$v$ corpus, pass each through the trained model, and extract the
projected token embedding at $w$'s position. The word-level embedding is the
mean over all occurrences:

\begin{equation}
    \bE_v[w] = \frac{1}{|C_v(w)|} \sum_{c \in C_v(w)}
    \text{proj}(\text{BETO}_{\text{LoRA}}(c))[w]
\end{equation}

\subsection{Procrustes Alignment}

To bring all variety embeddings into a shared space, we apply orthogonal
Procrustes alignment \citep{schonemann1966generalized} using anchor words
(high-frequency words with minimal dialectal variation):

\begin{equation}
    \bE_v^{\text{aligned}} = \bE_v \mathbf{R}_v^*, \quad
    \mathbf{R}_v^* = \arg\min_{\mathbf{R}^\top \mathbf{R} = \bI}
    \|\bE_v[\text{anchors}] \mathbf{R} - \bE_{\text{ref}}[\text{anchors}]\|_F
\end{equation}

with ES\_PEN (Peninsular) as the reference variety.

\subsection{Transformation Matrices}

For each non-reference variety $v$, we compute the transformation matrix
$\bW_v$ via regularized least squares:

\begin{equation}
    \bW_v = \bE_v^{\text{aligned}} \bE_{\text{ref}}^{+}
\end{equation}

where $\bE_{\text{ref}}^{+}$ is the regularized pseudoinverse. This matrix
captures how variety $v$ transforms the reference embedding space.

\subsection{Eigendecomposition and Dialect Fingerprints}

Each transformation matrix is eigendecomposed:

\begin{equation}
    \bW_v = \bP_v \boldsymbol{\Lambda}_v \bP_v^{-1}
\end{equation}

The eigenvalue magnitudes $|\lambda_1| \geq |\lambda_2| \geq \cdots$ form
a \emph{spectral fingerprint} of each dialect. The spectral gap
$|\lambda_1| - |\lambda_2|$, effective rank (number of eigenvalues capturing
90\% of total magnitude), and the eigenvector directions all encode
dialect-specific information.

\subsection{DIAL: Dialectal Interpolation via Algebraic Linearization}

The eigendecomposition enables continuous intensity control:

\begin{equation}
    \bW_v(\alpha) = \bP_v \boldsymbol{\Lambda}_v^{\alpha} \bP_v^{-1}
\end{equation}

where $\alpha = 0$ gives the identity (reference dialect), $\alpha = 1$ gives
the full dialect transformation, and $\alpha > 1$ produces ``hyperdialect''
(exaggerated dialectal features).

\subsection{Scoring}

Given a new text, we:
\begin{enumerate}
    \item Tokenize and extract bag-of-words embedding as the mean of
    vocabulary-matched word embeddings.
    \item For each variety $v$, project onto the top eigenmodes of $\bW_v$ and
    compute a dialect affinity score based on alignment with
    variety-specific eigenmode directions.
    \item Normalize scores to probabilities via softmax.
\end{enumerate}

\section{Classification Methods}

We evaluate three classification approaches, all built on the same trained
model:

\subsection{ArcFace Direct Classification}

The ArcFace head produces scaled cosine logits directly. At inference
(without angular margin), the prediction is $\arg\max_j \hat{\mathbf{z}}^\top
\hat{\mathbf{w}}_j$.

\subsection{Projection-Space Centroid Classifier}

We compute per-variety centroids from corpus texts:

\begin{equation}
    \boldsymbol{\mu}_v = \frac{1}{|S_v|} \sum_{t \in S_v}
    \hat{\mathbf{z}}(t), \quad
    \hat{\boldsymbol{\mu}}_v = \boldsymbol{\mu}_v / \|\boldsymbol{\mu}_v\|
\end{equation}

Classification uses temperature-scaled cosine similarity to centroids:

\begin{equation}
    P(v \mid t) = \frac{\exp(\hat{\mathbf{z}}(t)^\top
    \hat{\boldsymbol{\mu}}_v / \tau_{\text{inf}})}
    {\sum_{v'} \exp(\hat{\mathbf{z}}(t)^\top
    \hat{\boldsymbol{\mu}}_{v'} / \tau_{\text{inf}})}
\end{equation}

with $\tau_{\text{inf}} = 0.1$. This method requires no additional training
beyond computing centroids from a small corpus sample (200 documents per
variety).

\subsection{Eigenmode Spectral Classifier}

Uses the spectral fingerprinting approach described in
Section~5. This method is more computationally expensive but provides
interpretable eigenmodes that correspond to linguistic dimensions of variation.

\section{Preliminary Results}

\begin{table}[h]
\centering
\caption{Classification accuracy on 12 dialectally-marked test sentences.}
\label{tab:results}
\begin{tabular}{lcc}
\toprule
\textbf{Method} & \textbf{Accuracy} & \textbf{Notes} \\
\midrule
Random baseline & 12.5\% & 1/8 varieties \\
Eigenmode spectral & 47.4\% (9/19) & 500 docs/variety sample \\
Projection centroid & 58.3\% (7/12) & 200 docs/variety centroids \\
ArcFace direct & 0\% & Collapsed (pre-fix) \\
\bottomrule
\end{tabular}
\end{table}

\begin{table}[h]
\centering
\caption{Per-variety projection centroid accuracy on test sentences.}
\label{tab:per_variety}
\begin{tabular}{lcl}
\toprule
\textbf{Variety} & \textbf{Correct/Total} & \textbf{Typical Confusion} \\
\midrule
ES\_RIO (Rioplatense) & 2/2 & --- \\
ES\_CHI (Chilean) & 2/2 & --- \\
ES\_MEX (Mexican) & 1/2 & $\to$ ES\_RIO \\
ES\_CAR (Caribbean) & 1/1 & --- \\
ES\_PEN (Peninsular) & 1/2 & $\to$ ES\_CAN \\
ES\_AND (Andalusian) & 0/1 & $\to$ ES\_PEN \\
ES\_CAN (Canarian) & 0/1 & $\to$ ES\_CAR \\
ES\_AND\_BO (Andean) & 0/1 & $\to$ ES\_RIO \\
\bottomrule
\end{tabular}
\end{table}

The projection centroid classifier shows strong performance on lexically
distinctive varieties (Rioplatense \emph{che, boludo, pibe}; Chilean
\emph{huev\'on, cachai, weá}; Caribbean \emph{chamo, vaina}) where
dialect-specific vocabulary provides clear signal. Performance is weaker on
Iberian varieties (Andalusian, Canarian, Peninsular) that share more lexical
overlap and differ primarily in phonological features that are partially
neutralized in written text.

\paragraph{Training dynamics.}
The contrastive loss plateaued at $\sim 2.65$ after 162K optimizer steps
with in-batch-only negatives (24 per query). This motivates the MoCo
integration, which provides 4096 consistent negatives and should enable
further contrastive descent.

\section{Discussion}

\paragraph{Feature space alignment matters.}
The ArcFace collapse demonstrates that multi-task objectives operating in
different feature spaces can interfere destructively. When ArcFace operated
on raw 768-dim pooled features while contrastive learning shaped the 384-dim
projected space, the pooled representation evolved to serve the projection
head but became degenerate for direct angular classification. Routing both
objectives through the same projected space resolves this.

\paragraph{MoCo vs.\ XBM.}
Cross-batch memory (XBM) stores training-encoder embeddings that become stale
as the model updates. With low temperature ($\tau = 0.07$), this staleness
is catastrophic: the contrastive loss converges to $\ln(Q) \approx 8.3$
(random baseline). MoCo's momentum encoder changes slowly enough
($m = 0.999$, so the half-life is $\sim 693$ steps) that queue entries
remain mutually consistent.

\paragraph{Spectral interpretability.}
The eigenmode approach offers interpretability that neural classifiers lack.
Each eigenmode corresponds to a direction of maximal dialectal variation,
and the eigenvalue magnitude indicates the strength of that variation axis.
The effective rank (number of eigenvalues capturing 90\% of total magnitude)
ranges from 108--150 across varieties, suggesting that dialectal variation
is distributed across many dimensions rather than concentrated in a few.

\section{Conclusion}

EigenDialectos combines transformer-based contrastive learning with spectral
analysis to classify Spanish text into 8 regional varieties. The
projection-space centroid classifier achieves 58\% accuracy (chance 12.5\%)
on a pre-MoCo model, with particularly strong performance on lexically
distinctive dialects. The MoCo integration, now implemented and validated
with 18 unit tests, should improve contrastive quality and consequently
classification accuracy. The spectral eigenmode approach provides an
interpretable complement to the neural classifier, with the DIAL framework
enabling continuous intensity control over dialectal transformations.

\paragraph{Future work.}
Training with MoCo (4096-entry queue) should push the contrastive loss
below the current plateau. The fixed ArcFace architecture (operating on
projected features) should provide a working direct classifier. Larger
evaluation sets with native-speaker annotations are needed for rigorous
accuracy assessment.

\bibliographystyle{plainnat}

\begin{thebibliography}{20}

\bibitem[Ca\~nete et~al.(2020)]{canete2020spanish}
J.~Ca\~nete, G.~Chaperon, R.~Fuentes, J.~Ho, H.~Kang, and J.~P\'erez.
\newblock Spanish pre-trained {BERT} model and evaluation data.
\newblock In \emph{PML4DC at ICLR}, 2020.

\bibitem[Deng et~al.(2019)]{deng2019arcface}
J.~Deng, J.~Guo, N.~Xue, and S.~Zafeiriou.
\newblock {ArcFace}: Additive angular margin loss for deep face recognition.
\newblock In \emph{CVPR}, 2019.

\bibitem[He et~al.(2020)]{he2020momentum}
K.~He, H.~Fan, Y.~Wu, S.~Xie, and R.~Girshick.
\newblock Momentum contrast for unsupervised visual representation learning.
\newblock In \emph{CVPR}, 2020.

\bibitem[Hu et~al.(2022)]{hu2022lora}
E.~J. Hu, Y.~Shen, P.~Wallis, Z.~Allen-Zhu, Y.~Li, S.~Wang, L.~Wang, and
W.~Chen.
\newblock {LoRA}: Low-rank adaptation of large language models.
\newblock In \emph{ICLR}, 2022.

\bibitem[Khosla et~al.(2020)]{khosla2020supervised}
P.~Khosla, P.~Teterwak, C.~Wang, A.~Sabel, Y.~Tian, C.~Isola,
A.~Maschinot, C.~Liu, and D.~Krishnan.
\newblock Supervised contrastive learning.
\newblock In \emph{NeurIPS}, 2020.

\bibitem[Rangel et~al.(2016)]{rangel2016overview}
F.~Rangel, P.~Rosso, B.~Verhoeven, W.~Daelemans, M.~Potthast, and B.~Stein.
\newblock Overview of the 4th author profiling task at {PAN} 2016.
\newblock In \emph{CLEF}, 2016.

\bibitem[Sch\"onemann(1966)]{schonemann1966generalized}
P.~H. Sch\"onemann.
\newblock A generalized solution of the orthogonal {Procrustes} problem.
\newblock \emph{Psychometrika}, 31(1):1--10, 1966.

\bibitem[Wang et~al.(2020)]{wang2020cross}
X.~Wang, H.~Zhang, W.~Huang, and M.~R. Scott.
\newblock Cross-batch memory for embedding learning.
\newblock In \emph{CVPR}, 2020.

\bibitem[Wen et~al.(2016)]{wen2016discriminative}
Y.~Wen, K.~Zhang, Z.~Li, and Y.~Qiao.
\newblock A discriminative feature learning approach for deep face recognition.
\newblock In \emph{ECCV}, 2016.

\bibitem[Yeh et~al.(2022)]{yeh2022decoupled}
C.~Yeh, C.~Hong, Y.~Hsu, T.~Liu, Y.~Chen, and Y.~LeCun.
\newblock Decoupled contrastive learning.
\newblock In \emph{ECCV}, 2022.

\bibitem[Zampieri et~al.(2017)]{zampieri2017findings}
M.~Zampieri, S.~Malmasi, N.~Ljubesi\'c, P.~Nakov, A.~Ali, J.~Tiedemann,
Y.~Scherrer, and N.~Aepli.
\newblock Findings of the {VarDial} evaluation campaign 2017.
\newblock In \emph{VarDial Workshop}, 2017.

\end{thebibliography}

\end{document}
